\section{Discussion}
\label{sec:discussion}

\todo{Write after Sections~\ref{sec:model}--\ref{sec:validation} are
complete with real numbers; a discussion drafted against placeholders risks
asserting conclusions the eventual evidence does not support. Points this
section should address, each contingent on what the completed validation
sections actually show:}

\begin{itemize}
\item \textbf{Where validation is strong and where it is not.} State
plainly which programs and geographies have committed, reproducible
validation evidence (federal and state income tax against \taxsim{}; nine
\soi{} baseline-level backtests; OBBBA and select tax-expenditure provisions
against \jct{}) and which do not (benefit programs generally have no
committed cross-model or administrative-outturn comparison surfaced during
this survey; \todo{confirm this gap is real and not just unsurveyed --
check policyengine-us's per-program test YAML for benefit programs against
known administrative benchmarks such as SNAP participation or issuance
totals before asserting the gap}).
\item \textbf{Scope of the \tpc{} comparison, if any is added.} No \tpc{}
comparison evidence was found during outline drafting; if v1 adds one, this
section must state clearly whether it is new analysis performed for the
paper or an existing artifact this outline missed.
\item \textbf{The in-sample/out-of-sample distinction.} Several of the
\jct{} figures cited in Section~\ref{sec:validation-jct-cbo} are calibration
targets rather than independent checks; the discussion must not conflate
"the model was calibrated to match X" with "the model independently
predicts X."
\item \textbf{Versioning discipline.} Revisit Section~\ref{sec:versioning}'s
claim about what does and does not warrant a v2, in light of whatever
validation drift (or stability) the completed Section~\ref{sec:validation}
actually shows across the \policyengineus{} versions the paper's cited
artifacts span.
\item \textbf{Limitations of the population-view for validation.} If
Section~\ref{sec:data} draws on population-view or harness concepts from
the imputation-paper / populace-fit line of work, be explicit about which
validations in this paper concern rules (statutory computation) versus
which concern data (microdata representativeness) -- they are validated by
different evidence and should not be presented as one number.
\end{itemize}
